\section{Separated pins and transfer to boundary brackets}\label{sec:pins}

The reflection norm of a Laplace insertion tends to zero by
Proposition~\ref{prop:reflection-null}.  We will transfer this identity to
upper and lower height laws whose boundary traces bracket zero.  We do
this by cutting the pinned boundary into separated pieces and moving
those pieces to one side of the insertion.  The same prescribed spins
can represent different integer heights on different pieces.  We first
require the pinned sites within each piece to lie in one consecutive
pair of heights, then require these pairs to have the same integer
offset.  Moment bounds and multiplicative comparisons of pin
probabilities control the motion.  Local phase factors allow the
common-offset condition to be imposed after the pieces have been moved
back.

All geometric lengths in this section are physical lengths.  A window,
a gap, or an arrangement of separated sets is fixed before the mesh
tends to zero.

An increment function supported in a disk means a function of the height
differences on edges in that disk, allowing a fixed arbitrarily small
enlargement for connecting paths.  If $f$ is smooth, supported in the disk,
and $\int f=0$, the exact area pairing $H_\delta(f)$ is such a function:
its coefficients sum to zero, so paths in the enlarged disk express it
in terms of increments.  Throughout the section put
\begin{equation}\label{pins:insertion-data}
 X_\delta=H_\delta(\Delta\varphi),\qquad
 F_\delta=\prod_{j=1}^q F_{j,\delta},\qquad
 \varphi\in C_c^\infty(K_0),\qquad
 \prod_{j=1}^q\|F_{j,\delta}\|_\infty\le C_F,
\end{equation}
where $\varphi$ is real, $K_0\Subset D$ is a disk, and $F_{j,\delta}$ is
an increment function in a disk $K_j\Subset D$.  The spectator disks may
overlap each other, but have positive distance from a fixed neighborhood
of $K_0$.  The empty product is allowed.

\subsection{The boundary walk and the two brackets}

Let $\gamma:\R/L\Z\longrightarrow\partial D$ be the positively oriented
arclength parametrization.  In a fixed tubular neighborhood of
$\partial D$, let $\pi$ be normal projection followed by the coordinate
$\gamma^{-1}$, with values in $\R/L\Z$.  This map exists because the
boundary is a $C^2$ embedded compact curve.  Write $d_L$ for circular
distance in $\R/L\Z$.

\begin{lemma}[Inner boundary walk]\label{lem:boundary-walk}
There is a closed walk $S_{0,\delta}$ of adjacent triangular-lattice face
centers, with repeated vertices and edges allowed, whose vertex set is
exactly the boundary face layer $\partial F_\delta$.  Its parametrized
trace is within $O(\delta)$ of the boundary cycle in traversal order.
Every free interior site lies in a bounded face of this walk, with
winding number one, and every interaction path from a free interior site
to an exterior site meets its vertex set.

For small $\delta$, $\pi$ is defined on the walk, its degree is one,
and the following fixed-scale assertions hold.  For each $b>0$, the
number of subwalks with disjoint traversal-time interiors and projected
diameter at least $b$ is bounded independently of sufficiently small
$\delta$.  Every proper projected boundary arc is traversed from end
to end by a subwalk, up to $o(1)$ errors at its endpoints.  These
assertions do not require a bound on the length of $S_{0,\delta}$ or
on its number of crossings of a line.
\end{lemma}

\begin{proof}
Associate to each successive edge of the honeycomb boundary cycle its
incident interior face.  At a trivalent boundary vertex, the two such
faces are either the same, or meet across the third edge.  Joining their
centers therefore gives the required walk after constant steps are
removed.  This also proves the $O(\delta)$ position bound.  Every face
adjacent to the boundary appears.  An edge path leaving the enclosed
face set first passes through such a face, proving separation.

One can deform each boundary edge to the corresponding center path in
the union of its incident faces and the stars of its endpoints.  This
deformation never meets a nonboundary interior face center.  Consequently
the original cycle and the center walk have the same winding number
around every free interior center.  In particular those centers lie in
bounded complementary components of the walk.

Use the uniformly convergent parametrizations in
Definition~\ref{def:approximation}, and parametrize the center walk in
the corresponding traversal order.  The position error still tends to
zero uniformly.  Normal projection is uniformly continuous on a smaller
closed tubular collar, so the projected walks converge uniformly to the
corresponding degree-one parametrization of the boundary circle.  They
have degree one for fine mesh by homotopy in that circle.  For any fixed
$b$, uniform continuity of the limiting parametrization gives
$\omega(b)>0$ such that a projected excursion of diameter $b$ needs at
least $\omega(b)$ units of parameter time, once the uniform error is
smaller than $b/4$.  Disjoint excursion intervals therefore number at
most $1/\omega(b)$, after normalizing the parameter circle to length one.
Finally lift the projected walk to the real line.  Its net change in
one traversal is $L$.  Before a first visit to the right endpoint of
a lifted proper arc, take its last visit to the left endpoint.  The
intervening subwalk stays in that interval and crosses it.  Discrete
endpoints introduce only $o(1)$ errors.  This proves the final assertion.
\end{proof}

In the following, the graph of a walk means its image with all traversed
edges.  Removing a set of its vertices removes the incident edges too.
Additional lattice adjacencies could only merge the components used
below, so the specified graph convention is harmless but convenient.

Fix a finite collection $\mathcal L$ of horizontal levels that meet
$D$ and are regular values of the vertical coordinate on $\partial D$.
We use nonempty collections, so their boundary intersection set is
nonempty.  Their
intersections with $\partial D$ are finite: each is isolated by
transversality, and their set is compact.  Let $\mathcal C\subset\R/L\Z$
be the finite set of projected intersection coordinates.  Choose
$\epsilon>0$ so small that the arcs
$\{d_L(s,c)<2\epsilon\}$, $c\in\mathcal C$, are disjoint and every
complementary arc has nonempty interior.  These arcs specify fixed
exceptional windows.  The pins will be exact off the smaller windows
$d_L(\pi(\cdot),\mathcal C)<\epsilon$.

Remove the vertices of $S_{0,\delta}$ with
$d_L(\pi(v),\mathcal C)<\epsilon/2$.  Among the remaining graph
components retain those that contain a vertex with
$d_L(\pi(v),\mathcal C)\ge\epsilon$; denote their union by
$A_{\delta,\epsilon}$.  Thus every skeleton vertex outside the
$\epsilon$ windows belongs to $A_{\delta,\epsilon}$, while this set
avoids the $\epsilon/2$ windows.

\begin{lemma}[Bounded label count]\label{pins:label-count}
At fixed $\mathcal L$ and $\epsilon$, the number of connected
components used to form $A_{\delta,\epsilon}$ is bounded by a constant
$M(\mathcal L,\epsilon)$ independent of sufficiently small $\delta$.
\end{lemma}

\begin{proof}
A retained component contains a traversal segment beginning next to a
removed interval and reaching distance at least $\epsilon$ from
$\mathcal C$, without visiting a removed vertex.  Indeed use a maximal
portion of the original traversal outside the removed set that visits
a deep vertex.  Its projection travels a distance at least
$\epsilon/2-o(1)$.  For distinct graph components these portions have
disjoint time interiors, because a shared traversal vertex would join
the components.  Lemma~\ref{lem:boundary-walk} bounds their number.
The same proof, using any fixed smaller constants in the two thresholds,
covers rounding at endpoints.  It puts no bound on the number of short
excursions confined to the windows.
\end{proof}

\begin{definition}[Filled brackets]\label{def:brackets}
In the plane spin law impose $B=+$ on all of $S_{0,\delta}$ and $W=+$
on $A_{\delta,\epsilon}$.  Anchor the height to zero at one doubly pinned
vertex, and restrict it to the original interior.  Its law is the
\emph{upper filled bracket}, denoted $H_\delta^{+,\epsilon}$.  For the
\emph{lower filled bracket} $H_\delta^{-,\epsilon}$ interchange the
colors: impose $W=+$ on the whole skeleton and $B=+$ on
$A_{\delta,\epsilon}$, with the same zero anchor.
\end{definition}

These finite pin events have positive probability.  In a finite free
ambient graph the constant zero height is an admissible extension;
Proposition~\ref{prop:rare-pins}, or its deterministic enclosing-cut
comparison, keeps the finite event positive in the plane limit.  All
plane conditionings below can equivalently be obtained by first taking
the ambient-volume limit at a fixed mesh.  Local comparison gives this
limit, and division is only by a fixed positive finite-event probability
at that stage.

\begin{lemma}[Integer traces and order]\label{pins:bracket-order}
The bracket traces are actual integer heights satisfying
\begin{equation}\label{pins:trace-intervals}
 H_\delta^{+,\epsilon}\in\{0,1\},\qquad
 H_\delta^{-,\epsilon}\in\{-1,0\}
 \quad\hbox{on }S_{0,\delta},\qquad
 H_\delta^{\pm,\epsilon}=0\quad\hbox{on }A_{\delta,\epsilon}.
\end{equation}
Conditional on either trace, the free interior has the original weighted
height Gibbs law.  There is an ordered coupling
\begin{equation}\label{pins:height-sandwich}
 H_\delta^{-,\epsilon}\le h_\delta\le H_\delta^{+,\epsilon}
\end{equation}
on every interior face.  The bounded-trace smearing estimates of
Proposition~\ref{prop:smearing} apply to both brackets with constants
independent of the window sizes.
\end{lemma}

\begin{proof}
The integer inverse image of $B=+$ is the union of the pairs
$4t+\{0,1\}$, whereas that of $W=+$ is the union of
$4t+\{-1,0\}$.  A nearest-neighbor chain of constant first color cannot
leave its lifted pair, since different pairs have distance at least
three.  Connectedness of $S_{0,\delta}$ and its zero anchor therefore
give precisely \eqref{pins:trace-intervals}; second-color plus pins
select the zero endpoint of the pair.  This argument establishes the
integer branch, rather than merely its residue modulo four.

Lemma~\ref{lem:boundary-walk} separates every free site from exterior
interactions.  After conditioning on the skeleton heights, the
nearest-neighbor weights factor over the interior components and the
exterior; factors between fixed sites affect only the trace probability.
Lemma~\ref{lem:spin-lift} identifies each interior factor with the
weighted height law.  Applying Lemma~\ref{lem:height-order} to each
fixed trace, and then averaging the resulting kernels, proves
\eqref{pins:height-sandwich}.  All these traces have absolute value at
most one, so Proposition~\ref{prop:smearing} gives the last assertion.
In particular the upper mean is between zero and one, and the lower
mean between minus one and zero, by comparison with the flat laws.
\end{proof}

Our target is the identity
\begin{equation}\label{eq:bracket-insertion}
 \lim_{\delta\downarrow0}
 \E\left[H_\delta^{\pm,\epsilon}(\Delta\varphi)
                       F_\delta\right]=0
\end{equation}
for the local factors in \eqref{pins:insertion-data}, with the grid and
windows fixed.  We will prove it under the separating geometric
hypotheses stated in Proposition~\ref{prop:bracket-insertion}: two grid
levels must bracket the insertion disk and miss the spectators, and the
remaining pieces in the intervening strip must separate to its two
sides in an oblique projection.  This condition will let us move all the
other factors to one side of a reflection line while keeping the moved
pin sets separated.  The next constructions provide the
integer and moment control needed to carry out those moves.

\subsection{Opening gaps and checking the local lift}

Choose two levels $\ell_-<\ell_+$ from $\mathcal L$ strictly below and
above $K_0$, respectively, and missing every spectator disk with a
positive margin.  Let $\mathcal C_0\subset\mathcal C$ be their boundary
intersection coordinates.  Fix $0<a\ll\epsilon$, smaller also than all
interpoint distances in $\mathcal C_0$.  Delete skeleton sites satisfying
$d_L(\pi(v),\mathcal C_0)<a$.  Of the remaining graph components keep
only those reaching distance at least $2a$ from $\mathcal C_0$.  Delete
the first-color pins on the discarded sites as well.  All deletions
lie in $O(a)$ neighborhoods of $\gamma(\mathcal C_0)$; no second-color
pin is changed.  The first color is $B$ for the upper bracket and $W$
for the lower one; the proof below uses upper-bracket color names.

Group the retained first-color sites according to the successive
projected arcs between points of $\mathcal C_0$.  Let $P_i$ be the spin
pin event on group $i$, including its second-color pins, and put
\begin{equation}\label{pins:rare-normalization}
 P=\bigcap_{i=1}^m P_i,\qquad p_{i,\delta}=\Prob(P_i),\qquad
 J_\delta=\prod_{i=1}^m p_{i,\delta}.
\end{equation}
The number $m$ is fixed.  It bounds the number of groups, not the number
of first-color graph components.  Each group contains a doubly pinned
representative $v_i$: the degree-one traversal crosses the intervening
arc, including its nonempty exact-pin core.  Transversality at
$\ell_\pm$ implies that, after the deletions, every group has a
positive distance from the two levels.  Thin neighborhoods with endcaps
around its limiting compact arc give pairwise disjoint filled isolating
neighborhoods $U_i$, each a topological disk.  Choose also a smaller
filled core $V_i\Subset U_i$ containing all its pin sites, with room
in $U_i\setminus V_i$ for two nested collar cuts.  These neighborhoods
avoid all fixed test disks.  All clearances are positive at this fixed
$a$.

Let $C_i$ be the additional local check that all first-color pinned
sites in group $i$ use one and the same lifted pair.  It is an event
of the spins and relative increments in $V_i$: choose connecting paths
there from $v_i$ to all the pinned sites and compare their lifts.
The result is independent of the selected paths by
Lemma~\ref{lem:spin-lift}.  Define
\begin{equation}\label{pins:checked-event}
 P_{\mathrm{loc}}=P\cap\bigcap_{i=1}^m C_i.
\end{equation}
The checks are not included in the probabilities $p_{i,\delta}$.

\begin{lemma}[Marker checks]\label{pins:marker-checks}
For every fixed $a,\epsilon$ and separated placement of the groups,
\begin{equation}\label{pins:check-success}
 \Prob(P_{\mathrm{loc}}\mid P)\longrightarrow1.
\end{equation}
The assertion holds also after any fixed rigid translations or lattice
reflections of the groups, and along converging placements that retain
the specified positive clearances.
\end{lemma}

\begin{proof}
Consider one intervening projected arc.  Every retained graph component
touches the boundary of one of the deleted intervals: otherwise it would
be a component of the original connected walk with no neighbor outside.
More explicitly, a maximal unremoved traversal interval containing a
vertex at depth $2a$ starts or ends at depth $a+o(1)$.  It gives a
pinned chain from that endpoint to depth $2a-o(1)$, hence a crossing
of the projected marker at depth $3a/2$.  There is also a retained
component crossing the whole intervening arc.  The last-left-visit,
first-right-visit construction in Lemma~\ref{lem:boundary-walk} gives
such a chain; since $a$ is much smaller than the arc length, it crosses
both endpoint markers.  Thus every component reaches at least one
marker, and one component joins the two markers.

Fix an endpoint marker $s_*$ and its spatial point $z_* =\gamma(s_*)$.
All crossings of that marker, however numerous, occur in a ball
$B(z_*,r_\delta)$ with $r_\delta\to0$.  This follows from uniform
convergence into the normal collar; an edge crossing the projected
coordinate adds only an $O(\delta)$ error.  Choose a fixed outer radius
$\rho\ll a$ such that $B(z_*,2\rho)\subset U_i$ and the ball avoids
all second-color pins and equality edges.  Each of the relevant pinned
chains enters the inner ball and leaves the outer one, because its
endpoints have projected distances at least $a/2-o(1)$ from $s_*$.
At these short scales arclength and Euclidean distance along the
$C^2$ boundary are uniformly comparable.

First fix $r>0$, and then take $\delta$ small enough that $r_\delta<r$.
By Lemma~\ref{lem:buffer-circuits}, a stack of favorable good circuits
between radii $r$ and $\rho$ fails with probability at most
$C(r/\rho)^\alpha$.  Apply this estimate keeping all second-color
constraints first as equality edges $E_0$.  They avoid this annulus,
and every first-color pin in it is favorable.  Restoring the actual
plus labels costs at most $2^{M(\mathcal L,\epsilon)}$ by
Lemma~\ref{pins:label-count} and the component-label construction of
Section~\ref{sec:discrete}.  A good surrounding circuit joins all the
pinned chains through its constant first-color boundary layer.  Every
joined chain therefore has the same actual lifted pair.  This costs
one circuit event at the marker, regardless of the number of chains.

There are finitely many markers.  After the mesh limit send $r$ to
zero to make all their failure probabilities vanish.  The traversing
component ties together the two marker pairs of each group, proving
\eqref{pins:check-success}.  Translations and reflections preserve this
geometric argument; apply the buffer estimate afresh in their images.
Only spins and increments enter $C_i$, so no reflection invariance of
the particular cell augmentation is being assumed.
\end{proof}

\subsection{Moments after rare conditioning}

Proposition~\ref{prop:rare-pins} applies to the ordinary events $P_i$,
because their filled neighborhoods are disjoint and their total number
of connected second-color prescriptions is bounded.  It gives
\begin{equation}\label{pins:multiplicative}
 c J_\delta\le\Prob(P)\le C J_\delta.
\end{equation}
Constants depend on the fixed separated geometry, including $a$.
This is a multiplicative estimate at the scale of the pin probabilities.
An additive mixing error would give no bound after division by
$J_\delta$, which may be extremely small.

\begin{proposition}[Checked pin moments]\label{prop:pin-moments}
Consider any of the fixed separated arrangements above, including
arrangements enlarged by reflected copies.  Let $Y_\delta=H_\delta(f)$,
where $f$ is a real smooth zero-integral test supported in a disk having a
filled cut-search neighborhood disjoint from all $U_i$.  For every
$1\le p<\infty$,
\begin{equation}\label{eq:pin-moments}
 \E\bigl[|Y_\delta|^p\mid P_{\mathrm{loc}}\bigr]\le C_p
\end{equation}
for all sufficiently small $\delta$.  The constants are uniform on
compact families of such arrangements with positive clearances and
bounded smooth test norms.  The same holds for finite collections of
these tests and, by separating real and imaginary parts, for complex
tests.
\end{proposition}

\begin{proof}
Work first in a finite simply connected free ambient triangulation,
large enough to contain all neighborhoods and collars.  On
$P_{\mathrm{loc}}$, anchor $H(v_1)=0$.  The constrained height arrays
are exactly those satisfying, for integers $t_i$ with $t_1=0$,
\begin{equation}\label{pins:offset-constraints}
 \begin{aligned}
 H(v)&\in4t_i+I\quad &&\text{on all first-color pins in group }i,\\
 H(v)&=4t_i&&\text{on its second-color pins},
 \end{aligned}
 \qquad
 I=\begin{cases}\{0,1\},&\text{upper bracket},\\
                  \{-1,0\},&\text{lower bracket}.
    \end{cases}
\end{equation}
The forward implication is the definition of the checks.  Conversely
these constraints impose both colors at their prescribed sites and one
pair on each whole group.  Simple connectedness and the Gray-code lift
exclude additional winding constraints.  The array weight is exactly
the product of the original edge weights.

Partition these arrays by
$\kappa=(t_2\bmod2,\ldots,t_m\bmod2)$, so there are at most
$2^{m-1}$ classes.  If $H,H'$ belong to the same class, put
$s_i=(t_i+t_i')/2\in\Z$.  On a first-color pinned site, each of
$\lfloor(H+H')/2\rfloor$ and $\lceil(H+H')/2\rceil$ belongs to
$4s_i+I$; on a double pin both equal $4s_i$.  Thus both rounded
outputs satisfy \eqref{pins:offset-constraints}, although their parity
classes can differ from the input class.  The anchor stays zero.
The weighted midpoint verification in Section~\ref{sec:moments}
establishes the edge-weight premise of Theorem~\ref{thm:midpoint}.  What
remains is a positive central mass in every input class.

In the test neighborhood require one good cut surrounding the test,
with a fixed positive separation from its support.  In the collar of
each nonanchor obstacle require two nested good cuts of opposite signs
of the first color in $U_i\setminus V_i$, enclosing the entire check support.
All these searches have disjoint fixed buffers free of
equality edges and second-color labels.  Lemma~\ref{lem:buffer-circuits},
successive conditioning, and the bounded label cost give a joint
conditional probability at least $c_0>0$ under $P$.  Expose the test
cut from the exterior and choose the collar cuts by their corresponding
stopped explorations.  Lemma~\ref{pins:marker-checks} implies that
intersecting this event with $P_{\mathrm{loc}}$ still has probability
at least $c_0/2$ under $P$ for fine mesh.

The pin checks can be read within the $U_i$, outside the entire test
search region.  Conditional on the exterior record and on the actual
height layer of the test cut, its inside is the unconstrained weighted
height problem with a consecutive-pair trace.  Subtract its possibly
random common offset.  Since $\int f=0$, this does not change
$Y_\delta$.  Proposition~\ref{prop:smearing}, using a finite partition
into smaller smearing balls if necessary, gives a fixed $M_0$ and a
conditional probability at least $1/2$ of $|Y_\delta|\le M_0$.
Consequently a central event $\mathcal A$, including all selected cuts
and checks and the condition $|Y_\delta|\le M_0$, has
\begin{equation}\label{pins:central-unpartitioned}
 \Prob(\mathcal A\mid P_{\mathrm{loc}})\ge c_1>0,
 \qquad \mathcal A\subset\{|Y_\delta|\le M_0\}.
\end{equation}

Lemma~\ref{lem:collar-switch} supplies in each nonanchor collar a
weight-preserving involution shifting all its inner heights by $+4$
or $-4$, relative to the exterior.  The relevant component is chosen
from the first-color, status, and cut record, without using its fair
second-color label.  Its flip preserves that record.  All spins on
the protected inner and outer regions, all pin labels, and all local
checks remain unchanged.  It also preserves $Y_\delta$ and its test-cut
record: the entire test search neighborhood lies outside $U_i$, in the
protected exterior of the switch.  The disjoint collars give commuting involutions;
the $i$th one flips precisely the parity of $t_i$.  These involutions
act on the augmented spin and status measure.  The event $\mathcal A$
records only the selected cuts, the checks, and the displayed bound on
$Y_\delta$; it does not fix the extra exterior heights used to estimate
its probability.  It is therefore invariant under the involutions.
Each orbit has one equally weighted
member for every parity vector, so
\begin{equation}\label{pins:every-class-mass}
 \Prob\bigl(|Y_\delta|\le M_0,\ \kappa=k
                   \mid P_{\mathrm{loc}}\bigr)
 \ge 2^{-(m-1)}c_1\quad\text{for every }k.
\end{equation}
Marginalizing the statuses preserves this identity for the height
parity classes.  The orbit calculation does not assert independence
after conditioning on the checks.  When $m=1$ no collar switch is needed.

The face-area weights of $f$ have bounded total absolute mass, and
coordinatewise rounding changes the associated linear statistic by
at most half that mass.  Increase $M_0$ to absorb this bound.  Apply
Lemma~\ref{lem:midpoint-amplification} with
\eqref{pins:every-class-mass} and the parity partition.  It gives
\eqref{eq:pin-moments}, uniformly in the ambient volume.  Exhaustion
at each fixed mesh gives the plane assertion; Fatou's lemma transfers
the moment bounds.  With no obstacles use
Proposition~\ref{prop:plane-moments}.  The construction uses only fixed
positive margins and finitely many requests, proving the stated
uniformity and the finite-collection version by H\"older's inequality.
\end{proof}

\subsection{Normalized reflection bounds and analytic translation}

A height difference between representatives of two obstacles is not
local to either obstacle.  Its exponential can be split into two local
phase factors, one at each representative.  We use the stationary
Haar-lift phases from Proposition~\ref{prop:reflection-transfer}: for any
fixed real frequency $\theta$, put
\begin{equation}\label{pins:haar-phase}
 Z_\theta(v)=U_\theta\exp\{i\theta(H(v)-H(o))\},
\end{equation}
with an independent uniform unit complex variable for every phase field
needed.  Two factors from the same field share this auxiliary variable,
so
$Z_\theta(v)\overline{Z_\theta(w)}
=\exp\{i\theta(H(v)-H(w))\}$.
Changing the anchor is absorbed by Haar invariance.  The joint spin,
increment, and phase law is reflection positive, and a phase value at a
site on one side of a row belongs to that side's algebra.  Thus each
factor can move with its own obstacle, including when its balancing
factor lies on the other side of the row.

For a site-row reflection $\vartheta$, use
$\langle A,B\rangle_\vartheta=\E[\overline{\vartheta A}B]$ on its
positive side, quotienting by null vectors.  Pure normal translation
by $b_\delta=\sqrt3\,\delta$ is the positive self-adjoint contraction
$S_\delta$ of Proposition~\ref{prop:reflection-transfer}.

\begin{lemma}[Normalized side norm]\label{pins:side-norm}
Suppose a family of checked obstacles and factors lies strictly on one
side of a site row.  Their filled neighborhoods and all needed check
paths lie on that side; unbounded factors have free test neighborhoods
disjoint from the obstacle neighborhoods.  If $A$ is a bounded product,
possibly including arbitrary Haar phases, times at most one smooth
zero-total insertion, then
\begin{equation}\label{pins:normalized-side-bound}
 \left\|A\prod_i\frac{\ind_{P_i\cap C_i}}{p_{i,\delta}}
       \right\|_\vartheta\le C.
\end{equation}
The bound is uniform on compact families with the stated margins.
\end{lemma}

\begin{proof}
Let $J$ denote the product of the probabilities on this side.  Squaring
the norm gives
\[
 J^{-2}\E\left[\overline{\vartheta A}A\,
    \ind_{P_{\mathrm{loc}}\cap\vartheta P_{\mathrm{loc}}}\right].
\]
The reflected arrangement consists of two copies of every obstacle,
all with disjoint filled neighborhoods.  The strict row margin
separates the two copies.  Their individual probabilities are the
same $p_{i,\delta}$, so \eqref{pins:multiplicative} bounds the
probability of the ordinary doubled pin event by $CJ^2$.
The checked event has no larger probability.  If $A$ contains an
insertion $Y$, Proposition~\ref{prop:pin-moments} in this doubled
arrangement and Cauchy--Schwarz bound
$\E[|Y\,\vartheta Y|\mid P_{\mathrm{loc}}\cap
\vartheta P_{\mathrm{loc}}]$ by a constant.  Bounded factors and
unit-modulus phases contribute only their uniform bounds.  The
factor $J^2$ cancels exactly, proving the assertion.
\end{proof}

For a related separating-line continuation, see
\cite[Section~7.1, Lemma~7.4]{DCKLM2026}.
Here we must also normalize by rare pin probabilities:
Lemma~\ref{pins:side-norm} supplies the bounded matrix elements before
the continuation is applied.  The following argument proves the exact
version needed for these checked obstacles.

\begin{lemma}[Analytic translation]\label{lem:analytic-translation}
Partition a finite arrangement of the preceding checked obstacles and
local factors into two families separated strictly in projection on a
lattice reflection normal $n$.  At most one factor is an unbounded
zero-total insertion.  Translate the high family by $sn$, and let
$M_\delta(s)$ be its plane expectation normalized by the product
$J_\delta$ of all individual ordinary pin probabilities, at lattice
normal-step translations.  If
\[
 M_\delta(s_\delta)\longrightarrow0
\]
for every $s$ in a nonempty open interval $I\subset(0,\infty)$ and
every lattice rounding $s_\delta\to s$, then $M_\delta(0)\to0$.
The same is true for a low family translated by $-sn$, and along any
converging sequence of other placements preserving the required margins.
\end{lemma}

\begin{proof}
Choose $\eta>0$ smaller than the actual projection gap between the
enlarged supports.  Back the high family up by
$\eta_\delta=N_\delta b_\delta\to\eta$, where $N_\delta$ is a
positive integer.  There is still a fixed gap; choose a site row in
it, converging to a line in the limiting gap.  Let $B_\delta$ be the
normalized high-side factor in this backed arrangement and let
$A_\delta$ be the reflected conjugate of the normalized low-side
factor.  Lemma~\ref{pins:side-norm} gives
$\|A_\delta\|_\vartheta,\|B_\delta\|_\vartheta\le C$.
Individual $p_{i,\delta}$ are invariant under lattice translation
and reflection, so their product is the same throughout the motion.

For $\Re z>-\eta_\delta$ set
\begin{equation}\label{pins:analytic-family}
 f_\delta(z)=\int_{[0,1]}
 \lambda^{(z+\eta_\delta)/b_\delta}\,
             \gamma_{A_\delta,B_\delta}(\dd\lambda),
\end{equation}
where the finite cross measure is constructed in
Lemma~\ref{spectra:transfer}, with its circle variable integrated out.
For $\Re w>0$ use $\lambda^w=e^{w\log\lambda}$ on $\lambda>0$
and give it value zero at $\lambda=0$.
Corollary~\ref{spectra:analytic} proves that this is holomorphic and
bounded by $\|A_\delta\|_\vartheta\|B_\delta\|_\vartheta$.
Thus $|f_\delta|\le C^2$, and all fine
meshes have the common domain $\Re z>-\eta/2$, which contains zero
as an interior point.

At $s_\delta=m b_\delta\ge0$, the exponent in
\eqref{pins:analytic-family} is the positive integer $m+N_\delta$.
The transfer interpretation therefore gives the exact equality
$f_\delta(s_\delta)=M_\delta(s_\delta)$.  This never evaluates a
zeroth spectral power at the kernel.  Cauchy's estimates give a
uniform derivative bound on compact subsets of the common half-plane,
so converging lattice roundings sample the same limiting function.
Every sequence has a subsequence on which $f_\delta$ converges locally
uniformly to a bounded holomorphic function.  It vanishes on $I$ by
hypothesis, hence throughout the half-plane by the identity theorem,
and in particular at zero.  Subsequence extraction proves the claim
for the original sequence.

Rigid translation preserves distances inside each family, and moving
the families apart increases their projection gap.  The preliminary
backoff was smaller than that gap.  Thin filled neighborhoods and
check paths can therefore be kept on their own sides at every step;
their reflected copies have the same separation.  These observations
verify the side-norm hypotheses, including for convergent variations of
the other placements.  Reversing $n$ proves the low-family version.
\end{proof}

\subsection{The three-direction geometry and phase projection}

Here is the precise geometry needed for the insertion identity.  Put
\begin{equation}\label{pins:three-normals}
 e=(0,1),\qquad u=(\sqrt3/2,1/2),\qquad
 u'=(-\sqrt3/2,1/2).
\end{equation}
The levels $\ell_\pm$ chosen above strictly bracket $K_0$ and avoid
every spectator disk.  After the boundary intersection gaps are opened,
each boundary group and each spectator support lies wholly above,
below, or between these levels.  Call the outer families $T$ and $B$.
Require the middle family to split into $L$ and $R$ with, for some
$\kappa>0$,
\begin{equation}\label{pins:middle-geometry}
 \sup_{z\in L}z\cdot u<\inf_{z\in K_0}z\cdot u-\kappa,
 \qquad
 \inf_{z\in R}z\cdot u>\sup_{z\in K_0}z\cdot u+\kappa.
\end{equation}
Suprema over a family mean over all its enlarged supports.  Its
obstacle neighborhoods, check paths, and representative phases are
included in those supports.  Empty families are omitted.  Choose the
neighborhoods sufficiently thin to preserve these inequalities.
Overlapping spectator disks must belong to the same family and move
together.  This geometry is used later only with $K_0$ and the strip
very small compared with their distance to all remaining supports.

\begin{lemma}[A local choice of the grid]\label{pins:local-geometry}
Fix $z_0\in D$, $\overline{B(z_0,\rho)}\subset D$, and an integer
$q\ge0$.  One can choose a finite grid of regular horizontal levels,
depending only on $z_0,\rho,q$, with the following property.  If
$K_0\subset B(z_0,\rho/1000)$ and at most $q$ spectator disks of
radii less than $\rho/(1000(q+1))$ lie outside $B(z_0,\rho)$, two
grid levels give all the strict geometry in
\eqref{pins:middle-geometry}.  The same grid works for any finite
number of such collections of spectator disks.
\end{lemma}

\begin{proof}
Translate $z_0$ to zero for this geometric calculation.  In each
interval $(\rho/40,\rho/20)$ and $(-\rho/20,-\rho/40)$ choose
$q+1$ regular levels with mutual distances exceeding
$\rho/(80(q+1))$.  They exist because regular values are dense
and each candidate interval has length $\rho/40$.  Their union is
the grid.  A spectator's closed vertical projection has length less
than $2\rho/(1000(q+1))$, so it meets at most one level in either
candidate family.  At least one level in each family therefore
avoids all $q$ disks, with a positive margin.  These two levels
strictly bracket $K_0$.

Every boundary point and middle-strip spectator point is outside
$B(0,\rho)$ and has $|y|<\rho/20$.  Thus
$|x|>\sqrt{399/400}\,\rho$, and
$|(\sqrt3/2)x+y/2|>\rho/2$, with the sign of $x$.
A connected middle boundary arc or spectator disk cannot change this
sign.  It consequently belongs wholly to $L$ or $R$, whereas the
$u$-projections of $K_0$ have absolute value at most $\rho/1000$.
Opening small intersection gaps and choosing sufficiently thin filled
neighborhoods preserves these strict inequalities.  The construction
of the grid preceded the choice of the disks, proving the last assertion.
\end{proof}

\begin{lemma}[Separated insertion]\label{pins:three-moves}
At fixed positive gaps satisfying \eqref{pins:middle-geometry}, and for
any product $\Psi_\delta$ of local Haar-phase factors at the
representatives with fixed frequencies,
\begin{equation}\label{pins:separated-identity}
 J_\delta^{-1}\E\left[
 X_\delta F_\delta\Psi_\delta\ind_{P_{\mathrm{loc}}}\right]
 \longrightarrow0.
\end{equation}
This holds for every converging lattice-rounded sequence of placements
with the stated strict separations.
\end{lemma}

\begin{proof}
Figure~\ref{fig:pins-five-moves} records the successive separating
projections.  Keep $X_\delta$ fixed and move all other factors, including
each check and each phase with its obstacle.  First translate $T$ by $a_+e$ and
$B$ by $-a_-e$.  Each is an entire extreme side across an existing
horizontal gap.  Choose these two distances sufficiently large that
$T$ is strictly above $K_0$ and $B$ strictly below it in
$u$-projection.  With \eqref{pins:middle-geometry}, the families
$Q_+=T\cup R$ and $Q_-=B\cup L$ are then strictly above and below
$K_0$, respectively, in that projection.

Translate $Q_+$ by $bu$, with $b$ sufficiently large that it is
vertically above $K_0$ and below $K_0$ in $u'$-projection.  Both
requirements are possible because
$u\cdot e=1/2$ and $u\cdot u'=-1/2$.  Translate $Q_-$ by $-cu$,
with $c$ so large that it is above $K_0\cup Q_+$ in
$u'$-projection.  Each of these moves increases an existing $u$ gap.
Finally translate $Q_-$ by $du'$.  It is an entire extreme side
across a strict $u'$ gap, and $u'\cdot e=1/2$, so large $d$ puts
it vertically above $K_0$ as well.  At every move, within-family
distances stay fixed and between-family projection gaps increase.
Thus no filled obstacle neighborhood meets another obstacle or the
insertion neighborhood.  The required thin local paths remain available.

All distances can be chosen in open ranges, in the order
$a_+,a_-,b,c,d$.  The translations can be made exactly on the lattice.
If $a_1=(1,0)$ and $a_2=(1/2,\sqrt3/2)$ generate $\T$, then
\[
 \sqrt3 e=-a_1+2a_2,\qquad
 \sqrt3 u=a_1+a_2,\qquad
 \sqrt3 u'=-2a_1+a_2.
\]
Round each physical distance to a multiple of $b_\delta$.  The five
rounding errors are $O(\delta)$; all the finitely many strict margins
therefore survive.  Site reflection rows in these directions have
spacing $b_\delta/2$, so a row can also be chosen in every fixed gap.

At the final placement there is a horizontal site row separating
$X_\delta$ from every other factor.  Reflection Cauchy--Schwarz,
Lemma~\ref{pins:side-norm}, and
Proposition~\ref{prop:reflection-null} bound the absolute value of
the normalized expectation by $C\|X_\delta\|_\vartheta\to0$.
Adding phases does not change this norm of the phase-free insertion.

For every fixed admissible $a_+,a_-,b,c$, the last assertion holds
for $d$ in a nonempty open interval of sufficiently large values and
for every lattice rounding.  Lemma~\ref{lem:analytic-translation}
therefore removes the $du'$ move.  The resulting assertion holds for
every sufficiently large $c$ with earlier distances fixed, so the same
lemma removes the $-cu$ move.  Continue in the order $bu$, $-a_-e$,
$a_+e$.  At each stage later distances may depend on the already fixed
earlier ones; the open-interval assertion is exactly what the analytic
lemma needs.  This proves \eqref{pins:separated-identity} at the original
placement.  Empty families simply omit their moves.
\end{proof}

\begin{figure}[!hbp]
\centering
\begingroup
\colorlet{pinsplus}{blue!60!black}
\colorlet{pinsminus}{orange!65!black}
\tikzset{
  pins gap/.style={densely dashed,black!45,line width=.45pt},
  pins move/.style={-{Stealth[length=2mm]},line width=.8pt},
  pins title/.style={font=\small\bfseries},
  pins order/.style={font=\footnotesize,align=center},
  pins axis/.style={-{Stealth[length=1.5mm]},black!50,line width=.4pt}
}
\newcommand{\pinsgroup}[4]{%
  \begin{scope}[shift={(#1,#2)}]
    \draw[draw=#4,fill=#4!8,line width=.65pt]
      (0,0) ellipse[x radius=.34cm,y radius=.23cm];
    \fill[#4] (-.23,.03) circle[radius=.018cm];
    \fill[#4] (.20,.10) circle[radius=.018cm];
    \fill[#4] (.23,-.09) circle[radius=.018cm];
    \node[font=\small,inner sep=0pt] at (0,0) {$#3$};
  \end{scope}%
}
\newcommand{\pinskzero}{%
  \draw[draw=black!60,fill=black!6,line width=.65pt]
    (0,0) circle[radius=.23cm];
  \node[font=\footnotesize,inner sep=0pt] at (0,0) {$K_0$};
}
\newcommand{\pinsprojection}[6]{%
  \draw[pins axis] (-2.30,-1.95)--(2.30,-1.95)
    node[right,black,font=\scriptsize] {$\pi_{#1}$};
  \draw[#3,line width=2.3pt] (-1.75,-1.95)--(-1.05,-1.95);
  \draw[black!65,line width=2.3pt] (-.35,-1.95)--(.35,-1.95);
  \draw[#6,line width=2.3pt] (1.05,-1.95)--(1.75,-1.95);
  \node[pins order,anchor=north] at (-1.40,-2.03) {$#2$};
  \node[pins order,anchor=north] at (0,-2.03) {$#4$};
  \node[pins order,anchor=north] at (1.40,-2.03) {$#5$};
}
\begin{tikzpicture}[x=1cm,y=1cm,font=\footnotesize]
  \begin{scope}
    \node[pins title] at (0,2.15) {(a) Initial placement};
    \draw[pins gap] (-2.30,.46)--(2.30,.46)
      node[right,black] {$\ell_+$};
    \draw[pins gap] (-2.30,-.46)--(2.30,-.46)
      node[right,black] {$\ell_-$};
    \pinskzero
    \pinsgroup{0}{.95}{T}{pinsplus}
    \pinsgroup{0}{-.95}{B}{pinsminus}
    \pinsgroup{-1.70}{0}{L}{pinsminus}
    \pinsgroup{1.70}{0}{R}{pinsplus}
    \draw[pins move,pinsplus] (0,1.23)--(0,1.72)
      node[midway,right=3pt] {$(1)\ a_+e$};
    \draw[pins move,pinsminus] (0,-1.23)--(0,-1.72)
      node[midway,right=3pt] {$(2)\ {-a_-e}$};
    \pinsprojection{e}{B}{pinsminus}{K_0,L,R}{T}{pinsplus}
    \node[pins order,text width=6.1cm] at (0,-2.80)
      {$L<_uK_0<_uR$;\\the two horizontal gaps are strict.};
  \end{scope}

  \begin{scope}[shift={(7.15,0)}]
    \node[pins title] at (0,2.15) {(b) After moves 1--2};
    \draw[pins gap] (1.414,-1.45)--(-.260,1.45);
    \draw[pins gap] (.260,-1.45)--(-1.414,1.45);
    \pinskzero
    \pinsgroup{.90}{.95}{T}{pinsplus}
    \pinsgroup{1.60}{.05}{R}{pinsplus}
    \pinsgroup{-.90}{-.95}{B}{pinsminus}
    \pinsgroup{-1.60}{-.05}{L}{pinsminus}
    \node[text=pinsplus] at (1.60,1.45) {$Q_+=T\cup R$};
    \node[text=pinsminus] at (-1.60,-1.45) {$Q_-=B\cup L$};
    \draw[pins move,pinsplus] (2.00,.27)--(2.69,.67)
      node[midway,above=4pt] {$(3)\ bu$};
    \draw[pins move,pinsminus] (-2.00,-.27)--(-2.69,-.67)
      node[midway,below=4pt] {$(4)\ {-cu}$};
    \pinsprojection{u}{Q_-}{pinsminus}{K_0}{Q_+}{pinsplus}
    \node[pins order,text width=6.1cm] at (0,-2.80)
      {Move each entire family outwards\\across its $u$-gap.};
  \end{scope}

  \begin{scope}[shift={(0,-5.65)}]
    \node[pins title] at (0,2.15) {(c) After moves 3--4};
    \draw[pins gap] (-1.299,-1.45)--(.375,1.45);
    \draw[pins gap] (-.375,-1.45)--(1.299,1.45);
    \pinskzero
    \pinsgroup{1.80}{1.18}{T}{pinsplus}
    \pinsgroup{2.28}{.55}{R}{pinsplus}
    \pinsgroup{-1.80}{-1.18}{B}{pinsminus}
    \pinsgroup{-2.28}{-.55}{L}{pinsminus}
    \node[text=pinsplus] at (1.95,1.63) {$Q_+$};
    \node[text=pinsminus] at (-1.85,-1.63) {$Q_-$};
    \draw[pins move,pinsminus] (-2.38,-.17)--(-2.99,.18)
      node[midway,above=5pt] {$(5)\ du'$};
    \pinsprojection{u'}{Q_+}{pinsplus}{K_0}{Q_-}{pinsminus}
    \node[pins order,text width=6.1cm] at (0,-2.80)
      {$K_0<_eQ_+$;\\now $Q_-$ is the extreme high $u'$-side.};
  \end{scope}

  \begin{scope}[shift={(7.15,-5.65)}]
    \node[pins title] at (0,2.15) {(d) After move 5};
    \draw[pins gap] (-2.70,.40)--(2.70,.40)
      node[right,black] {$\vartheta$};
    \pinskzero
    \pinsgroup{1.30}{1.33}{T}{pinsplus}
    \pinsgroup{2.10}{.80}{R}{pinsplus}
    \pinsgroup{-1.30}{.80}{B}{pinsminus}
    \pinsgroup{-2.10}{1.33}{L}{pinsminus}
    \node[text=pinsplus] at (1.80,1.77) {$Q_+$};
    \node[text=pinsminus] at (-1.80,1.77) {$Q_-$};
    \node[anchor=north] at (0,-.29) {$X_\delta$ stays fixed in $K_0$};
    \begin{scope}[shift={(0,-1.47)}]
      \draw[pins move,black!70] (0,0)--(0,.56)
        node[above,black] {$e$};
      \draw[pins move,black!70] (0,0)--(.65,.375)
        node[right,black] {$u$};
      \draw[pins move,black!70] (0,0)--(-.65,.375)
        node[left,black] {$u'$};
    \end{scope}
    \draw[pins axis] (-2.30,-1.95)--(2.30,-1.95)
      node[right,black,font=\scriptsize] {$\pi_e$};
    \draw[black!65,line width=2.3pt] (-1.75,-1.95)--(-1.05,-1.95);
    \draw[pinsplus,line width=2.3pt] (.55,-1.92)--(1.75,-1.92);
    \draw[pinsminus,line width=2.3pt] (.55,-1.98)--(1.75,-1.98);
    \node[pins order,anchor=north] at (-1.40,-2.03) {$K_0$};
    \node[pins order,anchor=north] at (1.15,-2.03) {$Q_+\cup Q_-$};
    \node[pins order,text width=6.1cm] at (0,-2.80)
      {A horizontal reflection row\\separates the insertion.};
  \end{scope}
\end{tikzpicture}
\endgroup

\caption{The five moves in Lemma~\ref{pins:three-moves}.
The normal projection is $\pi_n(z)=z\cdot n$, and $A<_nB$ means
$\sup_{z\in A}z\cdot n<\inf_{z\in B}z\cdot n$ for the enlarged
supports.  The interval bars record this order; panels do not share a
common scale.  Each family carries its obstacle neighborhoods, checks,
local factors, and representative phases together.  Choose
$a_+,a_-,b,c,d$ successively in nonempty open ranges, each depending on
the earlier choices.  After the final reflection bound, analytic
translation removes the moves in reverse order $d,c,b,a_-,a_+$ by
successive continuation of the normalized expectation.  Empty families
omit their moves.}
\label{fig:pins-five-moves}
\end{figure}

On $P_{\mathrm{loc}}$ the double-pin representatives satisfy
$H(v_i)=4t_i$ in the anchor convention of
\eqref{pins:offset-constraints}.  In particular
\[
 D_i=\frac{H(v_i)-H(v_1)}4=t_i-t_1\in\Z.
\]
For $s=(s_2,\ldots,s_m)\in[0,2\pi]^{m-1}$, independent Haar
fields give the exact balanced character identity
\begin{equation}\label{pins:balanced-characters}
 \exp\left(i\sum_{i=2}^m s_i D_i\right)
 =\prod_{i=2}^m
 Z^{(i)}_{s_i/4}(v_i)\overline{Z^{(i)}_{s_i/4}(v_1)}.
\end{equation}
The auxiliary Haar variables cancel.  This identity is valid for
unbounded integer lifts; no reduction of the differences modulo four
has been made.  Let
\begin{equation}\label{pins:common-offset-event}
 P_= =P_{\mathrm{loc}}\cap\{D_i=0\text{ for all }i\ge2\}.
\end{equation}
Fourier orthogonality on the integer variables gives
\begin{equation}\label{pins:fourier-projection}
 \ind_{P_=}=
 \frac{\ind_{P_{\mathrm{loc}}}}{(2\pi)^{m-1}}
 \int_{[0,2\pi]^{m-1}}
       \exp\left(i\sum_{i=2}^m s_iD_i\right)\,\dd s.
\end{equation}
For each fixed $s$, Lemma~\ref{pins:three-moves} applies to the
right-hand character.  At this fixed geometry its normalized integrand
is bounded uniformly in $s$ and $\delta$, since
\[
 J_\delta^{-1}\E[|X_\delta F_\delta|\ind_{P_{\mathrm{loc}}}]
 \le \frac{\Prob(P)}{J_\delta}\,C_F
       \E[|X_\delta|\mid P_{\mathrm{loc}}]\le C.
\]
Equations~\eqref{pins:multiplicative} and \eqref{eq:pin-moments}
justify this bound.  Dominated convergence in
\eqref{pins:fourier-projection} yields
\begin{equation}\label{pins:common-offset-numerator}
 J_\delta^{-1}\E[X_\delta F_\delta\ind_{P_=}]
 \longrightarrow0.
\end{equation}

\subsection{Synchronizing offsets and filling the gaps}

Keep the grid and windows fixed.  Choose a radius $r_0>0$ so small
that the balls about $\gamma(\mathcal C_0)$ of radius $2r_0$ are
disjoint, avoid every test disk, and lie within the single-color portions
of the windows.  In particular they meet no second-color equality
edge or pin.  The first-color deletions are contained in balls of
radius $Ca$ about these points, for a fixed $C$.  Write
$\mu_{\delta,a}$ for the law conditioned on $P$,
$\mu_{\delta,a}^{=}$ for the law conditioned on $P_=$, and
$\mu_{\delta,0}$ for the corresponding filled bracket.  These denote
one chosen bracket sign at a time.

\begin{lemma}[Gap closing]\label{pins:gap-closing}
There is a function $\omega(a)\to0$ as $a\downarrow0$ such that
\begin{align}
 \limsup_{\delta\downarrow0}
       \Prob(P_=^c\mid P)&\le\omega(a),
       \label{pins:offset-synchronization}\\
 \limsup_{\delta\downarrow0}
 \left\|\mu_{\delta,a}^{=}|_{\mathcal I_\delta}
              -\mu_{\delta,0}|_{\mathcal I_\delta}\right\|_{\TV}
       &\le\omega(a),\label{pins:increment-tv}\\
 \limsup_{\delta\downarrow0}
       \E_{\mu_{\delta,a}^{=}}|X_\delta|^p
       &\le C_p\qquad(1\le p<\infty),
       \label{pins:uniform-common-moments}
\end{align}
where $\mathcal I_\delta$ is the joint increment information in the
fixed test disks.  The constants in the last line are independent of
all sufficiently small $a$.
\end{lemma}

\begin{proof}
For each whole gap use a logarithmic stack between radii $Ca$ and
$r_0$, retaining the second-color prescriptions first as equality
edges $E_0$.  All first-color pins in this stack are favorable and
there is no equality edge crossing it.  A surrounding good plus circuit
joins the traversing first-color arms on its two sides.  The arms
exist by the degree-one traversal in Lemma~\ref{lem:boundary-walk};
they run from the inner gap neighborhood to beyond the chosen outer
collar.  Together with the local checks, such a circuit identifies
the lifted pairs of the neighboring groups.  The stack failure
probability is at most $C(a/r_0)^\alpha$.  Restoring the boundedly
many second-color plus labels costs at most a fixed factor $2^M$.
There are finitely many gaps, and their successive connections join
all groups cyclically.  Lemma~\ref{pins:marker-checks} thus proves
\eqref{pins:offset-synchronization}, with a bound of this power form
plus terms whose mesh limsup is zero.

For comparison with the filled bracket, keep exactly the same
second-color equalities $E_0$ in both laws.  Add the missing favorable
first-color pins in one gap and use the ordered stopped coupling of
Proposition~\ref{prop:history-order}.  Search outward from the changed
region for a common good plus circuit in the same stack.  Except
with probability $C(a/r_0)^\alpha$, such a cut is found before the
search meets any equality edge or label prescription.  The only
changes are on its near side, so the correct stopped kernels agree
on its far side by Lemma~\ref{lem:cell-cut}.  Sample those remaining
variables identically.  This matches the spins and increments in all
test disks, and also the second-color data needed to read every
prescribed plus label.  The exterior of the small gap ball contains
paths within each test disk, so the matching of pair variables there
indeed matches its full local increment information.

Repeat for the finitely many gaps, including all discarded shallow
components inside the changed balls.  The conditional laws at every
stage use the common $E_0$ and their own correct query marginals.
After this coupling restore the label event.  Under each equality law
its probability is at least $2^{-M}$.  If the joint records match the
label data outside an exceptional event of probability $r$, then for
any event of the test data its two label-conditioned probabilities
differ by at most $C2^M r$.  This follows by comparing the two joint
numerators and the two label-event denominators separately.  Thus the
bounded label cost preserves the vanishing total-variation error.
Finally conditioning the gapped law further on $P_=$ changes its
total variation by at most $\Prob(P_=^c\mid P)$, which proves
\eqref{pins:increment-tv}.  No estimate here divides an additive
error by the probability of the rare first-color pin event: all
couplings already have those conditional laws as their marginals.

It remains to prove \eqref{pins:uniform-common-moments}; the separated
estimate \eqref{eq:pin-moments} alone does not provide a uniform
constant as the obstacle gaps close.  Fix one cut-search annulus
around the test support, independent of $a$, away from the boundary
collar.  Under $P$ it has a favorable good cut with conditional
probability at least $c>0$, uniformly in $a$: use the buffer bound
with all remote equalities retained and then the same bounded label
cost.  By \eqref{pins:offset-synchronization}, for all sufficiently
small $a$ its intersection with $P_=$ has conditional probability
at least $c/2$ in mesh limsup.  More precisely choose $a$ so that
$\omega(a)<c/4$, and then take the mesh sufficiently small.

The common-offset event can be read outside this whole search disk.
Indeed its local checks use the boundary neighborhoods, and one can
join all representatives by paths in the complement of the disk.
Expose the cut from outside.  Conditional on the exterior and its
height trace, the interior is again the unconditioned weighted
height problem with a consecutive-pair trace.  Subtracting its
offset and applying Proposition~\ref{prop:smearing} provides a
central interval $|X_\delta|\le M_1$ with positive conditional mass,
where $M_1$ is independent of $a$.

In the anchored array description of $P_=$ every offset is zero.
Thus the constraints are just the fixed intervals $I$ and fixed
singletons $\{0\}$ in \eqref{pins:offset-constraints}.  Both rounded
midpoints of any two such arrays are admissible, without a parity
partition.  Lemma~\ref{lem:midpoint-amplification}, with the preceding
uniform central mass and bounded rounding error, proves
\eqref{pins:uniform-common-moments}.  Ambient exhaustion is justified
as in Proposition~\ref{prop:pin-moments}.  The filled brackets have
the corresponding uniform moments by Lemma~\ref{pins:bracket-order}.
\end{proof}

\begin{proposition}[Bracket insertion identity]\label{prop:bracket-insertion}
Fix a grid, windows, and local factors as in
\eqref{pins:insertion-data}.  Suppose two grid levels strictly bracket
$K_0$, miss every spectator disk, and give the strict middle-strip
geometry \eqref{pins:middle-geometry} for the boundary groups after
arbitrarily small intersection gaps are opened.  Then both filled
brackets satisfy the target identity~\eqref{eq:bracket-insertion}.
In particular, any subsequential distributional bracket limit with
convergent local test statistics satisfies the corresponding identity
against products of bounded continuous functions of those statistics.
\end{proposition}

\begin{proof}
At each fixed sufficiently small $a$, equations
\eqref{pins:multiplicative} and \eqref{pins:offset-synchronization}
give $\liminf_{\delta\downarrow0}\Prob(P_=)/J_\delta>0$.
Divide \eqref{pins:common-offset-numerator} by this quantity to obtain
$\E_{\mu_{\delta,a}^{=}}[X_\delta F_\delta]\to0$.
Let $T_M(t)=\max(-M,\min(t,M))$.  The second moments in
\eqref{pins:uniform-common-moments} and in the filled bracket imply
\[
 \E[|X_\delta-T_M(X_\delta)|]\le C/M
\]
in both laws, uniformly in small $a$ after the mesh limit.  Total
variation in \eqref{pins:increment-tv} therefore yields
\[
 \limsup_{\delta\downarrow0}
       |\E_{\mu_{\delta,0}}[X_\delta F_\delta]|
 \le 2C_F M\omega(a)+2C_F C/M.
\]
First take the mesh limit at fixed $a$ as displayed, then let
$a\downarrow0$, and finally let $M\to\infty$.  This proves
\eqref{eq:bracket-insertion}.  Weak convergence of the local test
statistics, continuity of their bounded factors, and the same moment
bound pass the identity to subsequential bracket limits.

The window size $\epsilon$ has stayed fixed throughout.  The boundary
estimates in Section~\ref{sec:boundary} will control the effect of
shrinking these windows, and Section~\ref{sec:limit} will then transfer
the identity to the true field.
In particular no separated-obstacle constant was used uniformly when
the gaps collided; the independent common-offset estimate was the
ingredient that permitted their closure.
\end{proof}
